\section{Übung 7 – Zufallsereignisse und statistische Abhängigkeit}

Zwei bzw.\ sechs ideale Würfel werden geworfen. Alle Teilaufgaben folgen dem
Laplace-Modell (vgl.\ Abschnitt~\ref{sec:axiome-wahrscheinlichkeit}).

% ----------------------------------------------------------------
\subsection*{Teil A – Zwei Würfel}
% ----------------------------------------------------------------

Der Ergebnisraum ist das kartesische Produkt
$\Omega = \Omega_w \times \Omega_w$ mit $|\Omega| = 6^2 = 36$
(vgl.\ Abschnitt~\ref{sec:kartesische-produktmengen}).
Jedes Elementarereignis $(\xi_i,\xi_j)$ hat die Wahrscheinlichkeit
$P(\xi_i,\xi_j) = \tfrac{1}{36}$.

% ---- A1 --------------------------------------------------------
\subsubsection*{A1 – Genau ein Würfel zeigt eine 5}

\begin{beispiel}[Genau eine 5 bei zwei Würfeln]

\textbf{Gegeben:} $n = 2$ unabhängige ideale Würfel, $p = P(\text{5}) = \tfrac{1}{6}$.

\textbf{Gesucht:} $P(\text{genau eine 5})$

\textbf{Formel} (Binomialformel, Abschnitt~\ref{sec:kombinatorik}):

$$
P = \binom{n}{s}\,p^{s}\,(1-p)^{n-s}
\qquad n=2,\; s=1,\; p=\tfrac{1}{6}
$$

\textbf{Rechnung:}

$$
P = \binom{2}{1}\cdot\frac{1}{6}\cdot\frac{5}{6}
  = 2 \cdot \frac{5}{36}
  = \frac{10}{36}
$$

\textbf{Lösung:}

$$
\boxed{P(\text{genau eine 5}) = \frac{5}{18} \approx 0{,}278}
$$

\end{beispiel}

% ---- A2 --------------------------------------------------------
\subsubsection*{A2 – Mindestens ein Würfel zeigt eine 5}

\begin{beispiel}[Mindestens eine 5 bei zwei Würfeln]

\textbf{Gegeben:} wie A1.

\textbf{Gesucht:} $P(\text{mindestens eine 5})$

\textbf{Formel} (Komplementregel + statistische Unabhängigkeit,
Abschnitt~\ref{sec:unabhaengige-versuche}):

$$
P(\text{mindestens eine 5}) = 1 - P(\text{keine 5})
= 1 - \left(\frac{5}{6}\right)^{2}
$$

\textbf{Rechnung:}

$$
P(\text{keine 5}) = \frac{5}{6}\cdot\frac{5}{6} = \frac{25}{36}
\qquad\Rightarrow\qquad
P = 1 - \frac{25}{36} = \frac{11}{36}
$$

\textbf{Lösung:}

$$
\boxed{P(\text{mindestens eine 5}) = \frac{11}{36} \approx 0{,}306}
$$

\end{beispiel}

% ---- A3 --------------------------------------------------------
\subsubsection*{A3 – Beide Würfel zeigen die gleiche Zahl}

\begin{beispiel}[Beide Würfel zeigen die gleiche Zahl]

\textbf{Gegeben:} $|\Omega| = 36$.

\textbf{Gesucht:} $P(W_1 = W_2)$

\textbf{Formel} (Laplace-Experiment, Abschnitt~\ref{sec:axiome-wahrscheinlichkeit}):

$$
P(A) = \frac{|A|}{|\Omega|}
$$

\textbf{Rechnung:}

Günstige Ereignisse: $A = \{(1,1),(2,2),(3,3),(4,4),(5,5),(6,6)\}$, also $|A| = 6$.

$$
P = \frac{6}{36}
$$

\textbf{Lösung:}

$$
\boxed{P(W_1 = W_2) = \frac{1}{6} \approx 0{,}167}
$$

\end{beispiel}

% ---- A4 --------------------------------------------------------
\subsubsection*{A4 – Beide Würfel zeigen unterschiedliche Zahlen}

\begin{beispiel}[Beide Würfel zeigen unterschiedliche Zahlen]

\textbf{Gegeben:} Ergebnis aus A3.

\textbf{Gesucht:} $P(W_1 \neq W_2)$

\textbf{Formel} (Komplementregel):

$$
P(W_1 \neq W_2) = 1 - P(W_1 = W_2)
$$

\textbf{Rechnung:}

$$
P = 1 - \frac{1}{6} = \frac{5}{6}
$$

\textbf{Lösung:}

$$
\boxed{P(W_1 \neq W_2) = \frac{5}{6} \approx 0{,}833}
$$

\end{beispiel}

% ----------------------------------------------------------------
\subsection*{Teil B – Sechs Würfel}
% ----------------------------------------------------------------

Der Ergebnisraum erweitert sich auf
$\Omega = \Omega_w^{\times 6}$ mit $|\Omega| = 6^6 = 46\,656$.

% ---- B1 --------------------------------------------------------
\subsubsection*{B1 – Genau ein Würfel zeigt eine 5}

\begin{beispiel}[Genau eine 5 bei sechs Würfeln]

\textbf{Gegeben:} $n = 6$ unabhängige ideale Würfel, $p = \tfrac{1}{6}$.

\textbf{Gesucht:} $P(\text{genau eine 5})$

\textbf{Formel} (Binomialformel, Abschnitt~\ref{sec:kombinatorik}):

$$
P = \binom{6}{1}\cdot\left(\frac{1}{6}\right)^{1}\cdot\left(\frac{5}{6}\right)^{5}
$$

\textbf{Rechnung:}

$$
P = 6 \cdot \frac{1}{6} \cdot \frac{5^5}{6^5}
  = 1 \cdot \frac{3125}{7776}
$$

\textbf{Lösung:}

$$
\boxed{P(\text{genau eine 5}) = \frac{3125}{7776} \approx 0{,}402}
$$

\end{beispiel}

% ---- B2 --------------------------------------------------------
\subsubsection*{B2 – Mindestens ein Würfel zeigt eine 5}

\begin{beispiel}[Mindestens eine 5 bei sechs Würfeln]

\textbf{Gegeben:} $n = 6$, $p = \tfrac{1}{6}$.

\textbf{Gesucht:} $P(\text{mindestens eine 5})$

\textbf{Formel} (Komplementregel, Abschnitt~\ref{sec:unabhaengige-versuche}):

$$
P(\text{mindestens eine 5}) = 1 - \left(\frac{5}{6}\right)^{6}
$$

\textbf{Rechnung:}

$$
\left(\frac{5}{6}\right)^{6} = \frac{15\,625}{46\,656}
\qquad\Rightarrow\qquad
P = 1 - \frac{15\,625}{46\,656} = \frac{31\,031}{46\,656}
$$

\textbf{Lösung:}

$$
\boxed{P(\text{mindestens eine 5}) = \frac{31\,031}{46\,656} \approx 0{,}665}
$$

\end{beispiel}

% ---- B3 --------------------------------------------------------
\subsubsection*{B3 – Alle sechs Würfel zeigen die gleiche Zahl}

\begin{beispiel}[Alle sechs Würfel zeigen die gleiche Zahl]

\textbf{Gegeben:} $|\Omega| = 6^6 = 46\,656$.

\textbf{Gesucht:} $P(W_1 = W_2 = \cdots = W_6)$

\textbf{Formel} (Laplace-Experiment, Abschnitt~\ref{sec:axiome-wahrscheinlichkeit}):

$$
P(A) = \frac{|A|}{|\Omega|}
$$

\textbf{Rechnung:}

Günstige Ereignisse: $\{(1,\ldots,1),(2,\ldots,2),\ldots,(6,\ldots,6)\}$, also $|A| = 6$.

$$
P = \frac{6}{6^6} = \frac{6}{46\,656} = \frac{1}{7\,776}
$$

\textbf{Lösung:}

$$
\boxed{P(W_1 = \cdots = W_6) = \frac{1}{7\,776} \approx 1{,}29 \cdot 10^{-4}}
$$

\end{beispiel}

% ---- B4 --------------------------------------------------------
\subsubsection*{B4 – Alle sechs Würfel zeigen unterschiedliche Zahlen}

\begin{beispiel}[Alle sechs Würfel zeigen unterschiedliche Zahlen]

\textbf{Gegeben:} $|\Omega| = 6^6 = 46\,656$.

\textbf{Gesucht:} $P(\text{alle Würfel verschieden})$

\textbf{Formel} (Permutationen ohne Wiederholung, Abschnitt~\ref{sec:kombinatorik}):

$$
|A| = n! \qquad \text{(alle Anordnungen von } \{1,2,3,4,5,6\}\text{)}
$$

$$
P(A) = \frac{n!}{|\Omega|}
$$

\textbf{Rechnung:}

$$
|A| = 6! = 720
\qquad\Rightarrow\qquad
P = \frac{720}{46\,656} = \frac{5}{324}
$$

\textbf{Lösung:}

$$
\boxed{P(\text{alle verschieden}) = \frac{5}{324} \approx 0{,}0154}
$$

\end{beispiel}

% ----------------------------------------------------------------
\subsection*{Teil C – Bildpunkte eines Schwarz-Weiß-Bildes}
% ----------------------------------------------------------------

\subsubsection*{C1 – Elementarereignisse eines einzelnen Bildpunkts}

\begin{beispiel}[Ereignismenge eines Bildpunkts]

\textbf{Gegeben:} Ein digitalisiertes Schwarz-Weiß-Bild. Jeder Bildpunkt kann genau
einen von zwei Zuständen annehmen: schwarz~(S) oder weiß~(W). Beide Zustände treten
mit gleicher Wahrscheinlichkeit auf.

\textbf{Gesucht:} Menge der Elementarereignisse $\Omega_P$ für einen einzelnen Bildpunkt.

\textbf{Formel} (Definition Elementarereignis, Abschnitt~\ref{sec:mengen-ereignisse}):

$$
\Omega_P = \{\xi_1, \xi_2, \ldots\}
$$

\textbf{Rechnung:}

Da schwarz und weiß die einzigen möglichen Zustände sind und sich gegenseitig ausschließen:

$$
\Omega_P = \{S,\, W\}
$$

Da es sich um ein Laplace-Experiment handelt (Abschnitt~\ref{sec:axiome-wahrscheinlichkeit}):

$$
P(S) = P(W) = \frac{1}{|\Omega_P|} = \frac{1}{2}
$$

\textbf{Lösung:}

$$
\boxed{\Omega_P = \{S,\, W\}, \qquad P(S) = P(W) = \frac{1}{2}}
$$

\end{beispiel}

% ---- C2 --------------------------------------------------------
\subsubsection*{C2 – Verbundwahrscheinlichkeiten zweier unabhängiger Bildpunkte}

\begin{beispiel}[Gemeinsame Betrachtung zweier Bildpunkte]

\textbf{Gegeben:} Zwei benachbarte Bildpunkte mit
$P(S) = P(W) = \tfrac{1}{2}$; die Bildpunkte sind statistisch unabhängig.

\textbf{Gesucht:} Wahrscheinlichkeiten aller Elementarereignisse des kombinierten
Experiments.

\textbf{Formel} (Kartesisches Produkt + Verbundwahrscheinlichkeit bei Unabhängigkeit,
Abschnitte~\ref{sec:kartesische-produktmengen} und~\ref{sec:verbundwahrscheinlichkeiten}):

$$
\Omega = \Omega_P \times \Omega_P, \qquad
P(A_i, B_j) = P(A_i)\cdot P(B_j)
$$

\textbf{Rechnung:}

$$
\Omega = \{(S,S),\,(S,W),\,(W,S),\,(W,W)\}, \qquad |\Omega| = 4
$$

\begin{center}
\begin{tabular}{c|cc}
  & $B = S$ & $B = W$ \\
  \hline
  $A = S$ & $P(S,S) = \tfrac{1}{2}\cdot\tfrac{1}{2} = \tfrac{1}{4}$
           & $P(S,W) = \tfrac{1}{2}\cdot\tfrac{1}{2} = \tfrac{1}{4}$ \\[4pt]
  $A = W$ & $P(W,S) = \tfrac{1}{2}\cdot\tfrac{1}{2} = \tfrac{1}{4}$
           & $P(W,W) = \tfrac{1}{2}\cdot\tfrac{1}{2} = \tfrac{1}{4}$ \\
\end{tabular}
\end{center}

Probe: $\displaystyle\sum_{i,j} P(A_i,B_j) = 4\cdot\frac{1}{4} = 1$ \checkmark

\textbf{Lösung:}

$$
\boxed{P(S,S) = P(S,W) = P(W,S) = P(W,W) = \frac{1}{4}}
$$

Alle vier Kombinationen sind gleichwahrscheinlich, da die Bildpunkte statistisch
unabhängig sind und jeder Zustand mit $\tfrac{1}{2}$ auftritt.

\end{beispiel}

% ---- C3–C5 Neue Verbundwahrscheinlichkeiten -------------------

\medskip
Für die folgenden Teilaufgaben C3–C5 gelten die Verbundwahrscheinlichkeiten:

\begin{center}
\begin{tabular}{c|cc|c}
  & $B = S$ & $B = W$ & $P(A_i)$ \\
  \hline
  $A = S$ & $0{,}3$ & $0{,}2$ & \\
  $A = W$ & $0{,}2$ & $0{,}3$ & \\
  \hline
  $P(B_j)$ & & & $1$ \\
\end{tabular}
\end{center}

wobei $A$ den linken und $B$ den rechten Bildpunkt bezeichnet.

\subsubsection*{C3 – Wahrscheinlichkeit: linker Bildpunkt ist schwarz}

\begin{beispiel}[Randwahrscheinlichkeit des linken Bildpunkts]

\textbf{Gegeben:} $P(S,S) = 0{,}3$,\; $P(W,S) = 0{,}2$,\; $P(S,W) = 0{,}2$,\; $P(W,W) = 0{,}3$.

\textbf{Gesucht:} $P(A = S)$

\textbf{Formel} (Randwahrscheinlichkeit, Abschnitt~\ref{sec:verbundwahrscheinlichkeiten}):

$$
P(A_i) = \sum_{j} P(A_i, B_j)
$$

\textbf{Rechnung:}

$$
P(A = S) = P(S,S) + P(S,W) = 0{,}3 + 0{,}2 = 0{,}5
$$

\textbf{Lösung:}

$$
\boxed{P(A = S) = 0{,}5}
$$

\end{beispiel}

\subsubsection*{C4 – Wahrscheinlichkeit: rechter Bildpunkt ist weiß}

\begin{beispiel}[Randwahrscheinlichkeit des rechten Bildpunkts]

\textbf{Gegeben:} wie C3.

\textbf{Gesucht:} $P(B = W)$

\textbf{Formel} (Randwahrscheinlichkeit, Abschnitt~\ref{sec:verbundwahrscheinlichkeiten}):

$$
P(B_j) = \sum_{i} P(A_i, B_j)
$$

\textbf{Rechnung:}

$$
P(B = W) = P(S,W) + P(W,W) = 0{,}2 + 0{,}3 = 0{,}5
$$

\textbf{Lösung:}

$$
\boxed{P(B = W) = 0{,}5}
$$

\end{beispiel}

\subsubsection*{C5 – Statistische Unabhängigkeit der beiden Bildpunkte}

\begin{beispiel}[Unabhängigkeitstest]

\textbf{Gegeben:} Verbundwahrscheinlichkeiten wie C3; Randwahrscheinlichkeiten aus C3/C4.

\textbf{Gesucht:} Sind die beiden Bildpunkte statistisch unabhängig?

\textbf{Formel} (Bedingung für statistische Unabhängigkeit,
Abschnitt~\ref{sec:verbundwahrscheinlichkeiten}):

$$
P(A_i, B_j) \stackrel{!}{=} P(A_i)\cdot P(B_j) \quad \forall\, i, j
$$

\textbf{Rechnung:}

Zunächst alle Randwahrscheinlichkeiten:

$$
P(A=S) = 0{,}3 + 0{,}2 = 0{,}5, \qquad P(A=W) = 0{,}2 + 0{,}3 = 0{,}5
$$
$$
P(B=S) = 0{,}3 + 0{,}2 = 0{,}5, \qquad P(B=W) = 0{,}2 + 0{,}3 = 0{,}5
$$

Gegenbeispiel: Es genügt ein einziger Fall, bei dem die Bedingung verletzt ist:

$$
P(A=S)\cdot P(B=S) = 0{,}5 \cdot 0{,}5 = 0{,}25 \;\neq\; P(S,S) = 0{,}3
$$

\textbf{Lösung:}

$$
\boxed{\text{Die Bildpunkte sind statistisch \textbf{abhängig}.}}
$$

Obwohl beide Randwahrscheinlichkeiten identisch zu Teil~C2 sind ($P = 0{,}5$), weichen
die Verbundwahrscheinlichkeiten vom Produkt der Randwahrscheinlichkeiten ab. Benachbarte
Bildpunkte in einem realen Bild sind typischerweise statistisch abhängig, da ähnliche
Farben in räumlicher Nähe häufiger gemeinsam auftreten.

\end{beispiel}
